\chapter{Storage Management and\\ Garbage Collection}
\vspace{2cm}
%\vspace{1cm}

As mentioned in Chapter~10, a graph reducer requires the support of a storage
management system which allocates cells on request, and recovers garbage
cells for subsequent re-use. Storage management and garbage collection is a
subject on which there is a large literature. Cohen [1981] gives an excellent
survey with a comprehensive list of references.

The purpose of this chapter is to sketch the standard algorithms, to give an
assessment of their characteristics, and to make a brief survey of more
recently developed techniques.

\section{Criteria for Assessing a Storage Manager}

When considering a garbage collection technique it is helpful to keep in mind
the criteria against which it should be assessed. The main ones are:

\begin{numbered}
    \item \textit{What are its overheads} (in space and time)? All garbage collection
    systems consume resources, both in the form of per-cell extra storage
    requirements and in the CPU cycles taken to perform the collection.

    \item \textit{Does it support compaction?} If a storage manager repeatedly allocates,
    recovers and re-allocates variable-sized cells, the free storage tends to
    become fragmented into many small separate blocks. This can mean that
    a cell cannot be allocated because no free block is large enough, even
    though the total free storage is adequate. This phenomenon is known as
    \textit{storage fragmentation} [Knuth, 1976], and it can only be avoided by
    periodically \textit{compacting} all the cells together at one end of the address
    space, so as to produce a large contiguous free area from which to
    allocate new cells. Compaction also has a beneficial effect on virtual
    memory performance.

    \item \textit{How well does it support a sparsely used heap and virtual memory?} There
    has been much recent interest in very large \textit{persistent} heaps. The idea of
    persistence is that a functional operating system, for example, could
    incorporate a filing system as part of its data structures in a very large
    heap, rather than treating the filing system as something external to the
    program (an idea first implemented in Multics). In such a system, only a
    small fraction of the heap will be in active use at any time, and a virtual
    memory system is essential to cache the active portion in fast memory.

    \item \textit{Can it operate on a parallel machine, or in real time?} One of the
    attractions of functional languages is that they offer a natural way to
    exploit the power of parallel architectures (see Chapter~24), which
    requires storage managers that are capable of running on such a dis\-tributed system.

    Some garbage collection techniques require the computation to be
    stopped while garbage collection takes place, leading to an `embarrassing
    pause' during which the system appears to do nothing. This is unacceptable
    in real-time applications, and garbage collectors have been
    proposed which work in parallel with the useful computation. Such
    parallel collectors may also be suitable for parallel architectures.

    \item \textit{What is the effect of heap occupancy?} The performance of some
    algorithms drops sharply when the heap gets full.

    \item \textit{Can it recover cyclic structures?}
\end{numbered}

\noindent
These issues are all discussed by Cohen.

\section{A Sketch of the Standard Techniques}

There are several well-known garbage collection techniques. Among these
are \textit{mark-scan}, \textit{copying} and \textit{reference-counting} garbage collectors. In this
section we will give a brief sketch of the algorithms and their characteristics.
Cohen [1981] is the reference where no reference is given explicitly.

Mark-scan algorithms operate in two phases. First, all accessible cells are
\textit{marked} by traversing the entire accessible structure. Then a linear scan
through memory recovers all unmarked cells.

Copying algorithms work by copying the entire accessible structure from
one portion of the address space (\textit{from-space}) into another (\textit{to-space}),
thereby leaving all the garbage behind in from-space. Cells being copied into
to-space are placed contiguously, beginning at one end of the space, and
hence when copying is complete there is a contiguous area in to-space from
which new cells can be allocated. When to-space fills up with new cells, the
spaces are \textit{flipped} (i.e.\ to-space becomes from-space and vice versa) and the
process is repeated. The algorithm is surprisingly simple, and is well described
in Baker's classic paper [1978].

Reference-counting relies on keeping an extra field, called the \textit{reference-count}, in each cell. The reference-count field holds the number of references
to the cell (i.e.\ the number of pointers to the cell). This count is incremented
whenever a pointer is duplicated, and decremented whenever a pointer is
discarded. When the reference-count drops to zero the cell must be garbage,
since no other cells point to it.

Against our criteria, the techniques have the following characteristics.

\begin{numbered}
    \item \textit{Overheads.} A \textit{mark-scan} collector requires a \textit{mark bit} in each cell to
    indicate that the cell has been visited. In addition it appears at first that
    the mark phase will require an auxiliary \textit{stack} to guide its recursive
    tree-walk. Furthermore, the only bound on the size of this stack is the
    number of cells in the heap, though this bound would only be attained in
    pathological cases. This would be a heavy price to pay, but fortunately
    the Deutsch--Schorr--Waite pointer-reversing algorithm [Schorr and
    Waite, 1967] reduces the space overheads of the mark phase algorithm to
    a single bit per cell (in addition to the mark bit). This algorithm was
    explained in a different context in Section~11.6.1.

    \textit{Copying} collectors appear to have a 100\% space overhead, but in a
    virtual memory system the semi-space that is not in use will be paged out,
    so there is very little overhead in fast memory. Even during copying,
    activity only takes place at two sites in the target semi-space (to-space),
    so only two pages of to-space need to be paged in.

    \textit{Reference-counting} collectors require a reference-count field in every
    cell. In principle this field should be as wide as an address, since every cell
    in the heap could point to a single cell, but in practice reference-counts
    are almost always small. Hybrid systems have therefore been proposed,
    which have a limited-width reference-count field. When there are too
    many references to a cell and the reference-count field overflows it is set
    to a special value meaning `infinity', which is never decremented (so the
    cell is then irrecoverable). Cells irrecoverable by reference-counting are
    subsequently recovered by an occasional invocation of a mark-scan or
    copying collector.

    Reference-counting collectors are also somewhat less easy to use.
    Great care must be taken in the implementation never to duplicate a
    reference without incrementing the reference-count, though this is not,
    of course, a criticism of the adequacy of the algorithm itself. More
    seriously, many extra storage accesses are required to update the
    reference-counts.

    \item \textit{Compaction.} Compaction can be combined with the scanning phase of a
    \textit{mark-scan} collector. This is usually done using \textit{sliding compaction}, in
    which cells are slid down to one end of the address space, maintaining
    their address order. This means that cells which point to each other will
    not normally end up physically adjacent.

    A \textit{copying} collector is inherently compacting, since the cells are copied
    into a contiguous area in to-space. Furthermore, it is fairly easy to
    arrange that cells that point to one another get copied into physically
    adjacent locations, which significantly improves \textit{locality} and gives
    opportunities for extra-compact list representations (\textit{cdr-coding})
    [Baker, 1978]. This process is sometimes called \textit{linearizing} since linked
    lists get copied into a contiguous linear area of store, and it further
    reduces the storage overheads of a copying collector. The improvement
    in locality may also give improved paging performance in a virtual
    memory system.

    \textit{Reference-counting} does not inherently perform any compaction, but
    there is no reason why a compactor could not run concurrently with a
    reference-counting garbage collector.

    \item \textit{Sparsely used heap/virtual memory.} \textit{Mark-scan} and \textit{copying} collectors
    visit all accessible cells, not just those in immediate use. In contrast,
    \textit{reference-counting} collectors visit only cells in current use. For heaps in
    which only a small fraction of the accessible data is in active use, this
    represents a strong advantage for reference-counting.

    Without compaction the accessible cells get thinly spread through the
    address space, giving appalling paging behavior. The locality-improving
    possibilities of copying collectors (or reference-counting plus a copying
    compactor), mentioned above, thus make them preferable to sliding
    compaction.

    \item \textit{Parallel machines and real-time performance.} Since garbage collection
    began, researchers have tried to find ways to run garbage collection in
    parallel with useful computation, in an endeavor to eliminate the
    `embarrassing pause'. For \textit{mark-scan} collectors this may be achieved by
    arranging that garbage collection is performed by a process (or
    processor) \textit{in parallel with} useful computation. The algorithm is, of
    course, more complicated [Steele, 1975; Kung and Wong, 1977; Dijkstra
    \textit{et al.}, 1978].

    For \textit{copying} collectors, an ingenious scheme allows the copying
    process to take place incrementally, a fixed small amount being
    performed whenever a cell is allocated [Baker, 1978]. This scheme
    increases the overheads of the useful computation somewhat, in both
    time and space, and fails completely if to-space fills up before copying is
    completed.

    \textit{Reference-counting} collectors are inherently distributed in time, and
    hence need no modification for real-time performance.

    \item \textit{Effect of heap occupancy.} The performance of \textit{mark-scan} and \textit{copying}
    collectors degrades sharply as the heap gets full, since all the accessible
    data have to be visited in order to recover the few unused cells.
    \textit{Reference-counting}, on the other hand, is unaffected by heap occupancy.

    \item \textit{Cyclic structures.} \textit{Mark-scan} and \textit{copying} collectors have no problem
    with cyclic structures, but \textit{reference-counting} cannot recover them. The
    reason is that when a cell refers to itself it may have a non-zero reference-
        count even though it is not accessible from anywhere else (and hence is
    garbage). Recent developments in this area are discussed below.
    \end{numbered}

    \section{Developments in Reference-counting}

    The overheads of reference-counting, and its inability to recover cyclic
    structures, have often led to its dismissal as a garbage collection technique,
    except in specialized contexts. However, recent work has made progress
    towards alleviating these problems, and the inherently real-time and distributed nature of reference-counting is becoming increasingly attractive as
    parallel architectures gain in importance.

    \subsection{Reference-counting Garbage Collection of Cyclic Structures}

    Hughes [1982] has suggested an extension to the conventional reference-counting algorithm that would allow it to reclaim circular structures, based on
    previous work by Bobrow [1980].

    The key idea is simple and elegant. We regard the accessible data in the
    heap as a \textit{directed graph}, and divide this graph into its \textit{strongly connected
    components}. In this context we recall that

    \begin{numbered}
    \item a graph is \textit{strongly connected} if, for any two nodes \ml{A} and \ml{B}, there is a path
    from \ml{A} to \ml{B}, and vice versa;

    \item a \textit{strongly connected component} of a graph is a maximal strongly
    connected subgraph.
    \end{numbered}

    \noindent
    Now, it is clear that

    \begin{numbered}
    \item if one node of a strongly connected component is accessible, then all its
    nodes are (and vice versa);

    \item if we coalesce all the nodes in each strongly connected component, then
    the resulting \textit{derived graph} is acyclic.
    \end{numbered}

    \noindent
    But now, since the derived graph is acyclic, it is amenable to conventional
    reference-counting garbage collection; and when a node of the derived graph
    becomes unreferenced, all the nodes of the corresponding strongly connected
    component have become unreferenced.

    Hughes therefore suggests adding a second reference-count field to each
    node, which either contains the \textit{shared reference-count} for the strongly
    connected component of which the node is a part, or is used to point at the
    node which does hold the shared reference-count. He gives algorithms for
    incrementally maintaining the information about which components are
    strongly connected, and shows that they are rather cheap, except where a
    strongly connected component is broken up.

    It appears that this technique can successfully alleviate the `circular data
    structure problem', and thus allow exploitation of the other desirable
    characteristics of reference-counting. Aertes [1981] describes an essentially
    identical system. Brownbridge [1985] describes a different reference-counting technique, also claimed to be capable of recovering circular
    structures.

    Interestingly, all of these techniques will only work for implementations
    free of side-effects. In other words, they will work for implementations of
    functional languages, but not for Lisp (at least, not if the program uses
    RPLACs). Perhaps this is a new point in favor of functional languages!

    \subsection{One-bit Reference-counts}

    The logical extreme of the limited width reference-count idea is a one-bit
    reference-count field. This is not a new idea [Wise and Friedman, 1977].

    Recently, however, a number of researchers noticed that instead of storing
    a one-bit reference-count \textit{in the cell} it would be possible to store the reference-count \textit{in the pointer}. A single bit in each pointer identifies the pointer as being
    a \textit{unique reference} or a \textit{shared reference}. Cells are created with a unique
    reference to them; when a unique reference is duplicated, \textit{both} copies become
    shared references. When a unique reference is discarded the cell to which it
    points can be immediately re-used; when a shared reference is discarded no
    recovery is possible. Like all elegant ideas it is marvellously obvious in
    retrospect.

    The principal advantage of storing the reference-count in the pointer is that
    it completely eliminates the extra store accesses required to increment and
    decrement reference-counts.

    The benefits of even such a narrow reference-count are dramatic. Stoye \textit{et
        al.} [1984] report that up to \ml{70}\% of all garbage cells are recovered immediately
    they become unused in the SKIM SK combinator reduction machine. Further
    performance improvement is gained in the SKIM implementation because
    reclaimed cells are often re-used immediately, rather than being attached to
    the free list.

    \subsection{Hardware Support for Reference-counting}

    Much of the time overhead of reference-counting would be alleviated if
    hardware support were available. Wise [1985] describes hardware for a
    `smart' memory module, capable of detecting when one pointer is overwritten
    with another. When this occurs, the module sends a `decrement' message to
    the module which holds the cell pointed to by the old pointer, and an
    `increment' message to the module which holds the cell pointed to by the new
    pointer.

    \section{Shorting out Indirection Nodes}

    In Chapter~12 we discussed the introduction of \textit{indirection nodes} to preserve
    sharing when updating the root of a redex with its result. This is the only
    purpose of indirection nodes, and in all other respects they are a burden on
    the implementation, since they take up storage, and have to be `jumped over'
    when traversing the graph.

    It turns out that a rather simple modification to the garbage collector can
    `short out' all the indirection nodes in a graph, so they are no longer required.
    Consider a mark-scan collector. When it reaches an indirection node during
    the mark phase, it does not mark the indirection node. Instead it overwrites
    the pointer to the indirection node with the contents of the indirection,
    effectively shorting it out. Thus:

    \begin{center}
        {%
            \sffamily
            \small
            \begin{tikzpicture}
                \def\cellwidth{0.8cm}
                \def\cellheight{0.7cm}
                \def\toprow{0}
                \def\botrow{-2.4}

                % Top row: left node
                \node[draw, rectangle, minimum width=\cellwidth, minimum height=\cellheight] (tl1) at (0,\toprow) {};
                \node[draw, rectangle, minimum width=\cellwidth, minimum height=\cellheight, anchor=west] (tl2) at (tl1.east) {};
                \node[draw, rectangle, minimum width=\cellwidth, minimum height=\cellheight, anchor=west] (tl3) at (tl2.east) {};

                % Top row: indirection node
                \node[draw, rectangle, minimum width=\cellwidth, minimum height=\cellheight] (ti1) at (5,\toprow) {IND};
                \node[draw, rectangle, minimum width=\cellwidth, minimum height=\cellheight, anchor=west] (ti2) at (ti1.east) {};

                % Top row: right node
                \node[draw, rectangle, minimum width=\cellwidth, minimum height=\cellheight] (tr1) at (9.6,\toprow) {};
                \node[draw, rectangle, minimum width=\cellwidth, minimum height=\cellheight, anchor=west] (tr2) at (tr1.east) {};
                \node[draw, rectangle, minimum width=\cellwidth, minimum height=\cellheight, anchor=west] (tr3) at (tr2.east) {};

                % Top row: pointers
                \draw[-{Stealth[length=3mm,width=2mm]}, thick] (tl3.center) -- (ti1.west);
                \draw[-{Stealth[length=3mm,width=2mm]}, thick] (ti2.center) -- (tr1.west);

                % Label
                \node[anchor=west] at (0,-1.2) {becomes};

                % Bottom row: left node
                \node[draw, rectangle, minimum width=\cellwidth, minimum height=\cellheight] (bl1) at (0,\botrow) {};
                \node[draw, rectangle, minimum width=\cellwidth, minimum height=\cellheight, anchor=west] (bl2) at (bl1.east) {};
                \node[draw, rectangle, minimum width=\cellwidth, minimum height=\cellheight, anchor=west] (bl3) at (bl2.east) {};

                % Bottom row: indirection node
                \node[draw, rectangle, minimum width=\cellwidth, minimum height=\cellheight] (bi1) at (5,\botrow) {IND};
                \node[draw, rectangle, minimum width=\cellwidth, minimum height=\cellheight, anchor=west] (bi2) at (bi1.east) {};

                % Bottom row: right node
                \node[draw, rectangle, minimum width=\cellwidth, minimum height=\cellheight] (br1) at (9.6,\botrow) {};
                \node[draw, rectangle, minimum width=\cellwidth, minimum height=\cellheight, anchor=west] (br2) at (br1.east) {};
                \node[draw, rectangle, minimum width=\cellwidth, minimum height=\cellheight, anchor=west] (br3) at (br2.east) {};

                % Bottom row: pointers
                \draw[-{Stealth[length=3mm,width=2mm]}, thick] (bi2.center) -- (br1.west);
                \draw[-{Stealth[length=3mm,width=2mm]}, thick] (bl3.center) -- ++(0,-1.0) -| (br1.south);
            \end{tikzpicture}
        }%
    \end{center}

    Since all pointers to indirection nodes will be updated in this way it follows
    that the indirection nodes themselves will be unreferenced (and unmarked),
    so they can be collected with the rest of the garbage. Not only does this save
    store, but it also saves time when following the pointers that have been
    updated. A very similar technique will work for a copying collector.

    \section{Exploiting Cell Lifetimes}

    Another approach recently suggested by Lieberman and Hewitt [1983] is
    based on the observation that

    \begin{quote}
        \small
        The longer a cell has lived,\\
        the longer it is likely to live.
    \end{quote}

    \noindent
    Consider, for example, a heap which includes a filing system. Many files will
    be unused for long periods, while data structures that are currently being
    processed will have relatively short lifetimes. A conventional copying
    collector will copy the entire filing system each time it runs -- a very wasteful
    activity, since it is unlikely to recover any space from the inactive majority of
    the filing system.

    Hewitt and Lieberman therefore suggest dividing the address space into
    \textit{regions} of increasing age. Most pointers point backwards in time (that is, if
    they cross region boundaries, they will mostly point from younger regions to
    older ones). Where pointers point from an older region into a younger one
    (only update operations can cause this), they are constrained to go via an \textit{entry
        table} associated with the younger region.

    Now the youngest region can be garbage-collected \textit{independently}, using
    Baker's algorithm, so long as we preserve all cells referenced from its entry
    table. In general, any region can be garbage-collected without touching any
    older information. So all we have to do is to garbage-collect young regions
    (where garbage collection will be fruitful) more often than older ones (where
    it will be less fruitful, but eventually necessary).

    \section{Avoiding Garbage Collection}

    Another approach to garbage collection is to try to avoid it altogether. Wadler
    [1984] suggests a technique for compiling a certain class of functional program
    into a \textit{finite state machine} with a fixed number of registers and no heap. This, in
    effect, performs memory allocation in advance (rather as a conventional
    Pascal program has no problem with memory allocation). He calls his
    compiler the \textit{listless transformer.}

    The functional programs to which this technique is applicable are, not
    surprisingly, those that can be evaluated using \textit{bounded internal storage}. This
    includes, for example, functions that find the length of a list, add up a list,
    concatenate or merge two lists, or divide a list into two lists of odd and even
    elements. It excludes, however, functions that sort a list, append a list to
    itself, or work on tree-shaped data.

    Clearly the applicability of the method is limited, but where appropriate it
    is extremely effective, since the finite state machine can be made very fast.

    Wadler has a working implementation of his listless transformer, written in
    KRC.

    \section{Garbage Collection in Distributed Systems}

    Efforts to develop garbage collectors which work in parallel with useful
    computation have gained new impetus with the advent of parallel architectures, where the problem generalizes to many computation and garbage
    collection processes.

    Most work has been addressed to architectures with a single large address
    space (\textit{closely coupled} systems), for example Hudak [1983a and 1983b] and
    Ben-Ari [1984].

    Other efforts have been directed towards \textit{loosely coupled} systems in which
    the heap is distributed between a number of processing elements, and
    accessing a cell held by another processing element is recognized as a
    relatively expensive operation. Examples include Mohamed-Ali [1984] and
    Hughes [1985].

    \section*{References}

    \begin{references}

        \item Aerts, J.P.H. 1981. \textit{Implementing SASL without Garbage Collection.} EUT-Report
        81-WSK-05. Eindhoven Univ.\ of Technology. November.

        \item Baker, H. 1978. List processing in real time on a serial computer. \textit{Communications of
            the ACM.} Vol.~21, no.~4, pp.~280--94.

        \item Ben-Ari, M. 1984. Algorithms for on-the-fly garbage collection. \textit{ACM TOPLAS.} Vol.
        6, no.~3, pp.~333--44.

        \item Bobrow, D.G. 1980. Managing reentrant structures using reference-counts. \textit{ACM
            TOPLAS.} Vol.~2, no.~3, pp.~269--73.

        \item Brownbridge, D. 1985. Cyclic reference-counting for combinator machines. In
        \textit{Proceedings of the IFIP Conference on Functional Programming and Computer
            Architecture, Nancy.} September.

        \item Cohen, J. 1981. Garbage collection of linked data structures. \textit{ACM Computing Surveys.}
        Vol.~13, no.~3, pp.~341--67.

        \item Dijkstra, E.W., Lamport, L., Martin, A.J., Scholten, C.S., and Steffens, E.F.M.
        1978. On-the-fly garbage collection -- an exercise in cooperation. \textit{Communications
            of the ACM.} Vol.~21, no.~11, pp.~966--75.

        \item Hudak, P. 1983a. \textit{Distributed Graph Marking.} Research report 268, Computer Science
        Dept, Yale. January.

        \item Hudak, P. 1983b. Distributed task and memory management. In \textit{Proceedings of the
            ACM Symposium on Principles of Distributed Computing}, pp.~277--89. August.

        \item Hughes, R.J.M. 1982. \textit{Reference-counting with Circular Structures in Virtual Memory
            Applicative Systems.} Programming Research Group, Oxford.

        \item Hughes, R.J.M. 1985. A distributed garbage collection algorithm. In \textit{Proceedings of the
            IFIP Conf.\ on Functional Programming and Computer Architecture, Nancy}, pp.
        256--72. Jouannaud (editor). LNCS 201, Springer Verlag. September.

        \item Knuth, D. 1976. \textit{The Art of Computer Programming.} Vol.~I, Section 2.5. Addison
        Wesley.

        \item Kung, H.T., and Wong, S.W. 1977. \textit{An Efficient Parallel Garbage Collection System,
            and its Correctness Proof.} Dept of Comp.\ Sci., Carnegie-Mellon Univ.
        September.

        \item Lieberman, H., and Hewitt, C. 1983. A real time garbage collector based on the
        lifetimes of objects. \textit{Communications of the ACM.} Vol.~26, no.~6, pp.~419--29.

        \item Mohamed-Ali, K.A. 1984. Object oriented storage management and garbage
        collection in distributed processing systems. PhD Thesis, report TRITA-CS-8406.
        Royal Institute of Technology, Stockholm. December.

        \item Schorr, H., and Waite, W. 1967. An efficient machine independent procedure for
        garbage collection. \textit{Communications of the ACM.} Vol.~10, no.~8, pp.~501--6.

        \item Steele, G.L. 1975. Multiprocessing compactifying garbage collection. \textit{Communications of the ACM.} Vol.~18, no.~9, pp.~495--508.

        \item Stoye, W.R., Clarke, T.J.W., and Norman, A.C. 1984. Some practical methods for
        rapid combinator reduction. In \textit{Proceedings of the ACM Symposium on Lisp and
            Functional Programming, Austin}, pp.~159--66. August.

        \item Wadler, P. 1984. Listlessness is better than laziness: lazy evaluation and garbage
        collection at compile-time. In \textit{Proceedings of the ACM Symposium on Lisp and
            Functional Programming, Austin}, pp.~45--52. August.

        \item Wise, D. 1985. Design for a multiprocessing heap with on-board reference-counting.
        In \textit{Functional Programming and Computer Architecture, Nancy}, pp.~289--304.
        Jouannaud (editor). LNCS 201. Springer Verlag.

        \item Wise, D.S., and Friedman, D.P. 1977. The one-bit reference-count. \textit{BIT.} Vol.~17, no.
        3, pp.~351--9.

    \end{references}